\begin{savequote}
\includegraphics[width=4cm]{./images/pallinger-christoph-aass/Larynxtubus_32890_textwidth}

{\sffamily\tiny Ch. Pallinger}
\end{savequote}
\chapter
[\CO{[AIR]} Erweitertes Airwaymanagement]
{\CC{[AIR]}\\Erweitertes Airwaymanagement}
% \AasRsN{RS VIII}
\label{AIR}
\label{topic:airway}
\label{topic:airwaymanagement-erweitert}

\ChapterIntro
{
~~~
}{%
\begin{description}
  \item[Maintainer:] Sebastian Gabriel
  \item[Autoren:] Diverse
  \item[Reviewer:] Standardreviewprozess
  \item[Version:] {\DocVersion} ({\DocVersionInfo})
  \item[SHA1:] {\DocGitCommit}
\end{description}

{\TxTeilkapitelAASS}
}

%\SectionUnderConstruction{}
\clearpage

Für Personal, welches nicht sedieren darf, gilt: Wenn der
Patient wieder spontan atmet, dann bleibt der Larynxtubus in
situ, solange bis der Patient würgt oder Anzeichen zeigt,
sich den Tbus selbst entfernen zu wollen. 

Kinder: Cave Spasmus

\section{Einleitung}

\TextSlide{\TxBeschreibung}
{Nach der Kontrolle der Vitalfunktionen kann sich
herausstellen, dass der Patient \EE{kein Bewusstsein} und
\EE{keine normale Atmung} hat, d.\,h. \EE{kein Kreislauf
vorhanden} ist. In dieser Situation liegt die
Notfalldiagnose
\T{Atem-/Kreislaufstillstand} vor. Die einzig
mögliche Therapie ist die (sofortige)
\EE{Herz-Lungen-Wiederbelebung} (\EE{Reanimation}). Dabei
versucht der Helfer den Kreislauf des Patienten \enquote{von
außen} so gut wie möglich aufrecht zu
erhalten. Dies geschieht einerseits durch Beatmung
um dem Körper Sauerstoff zuzuführen und andererseits durch Herzdruckmassage um die
mechanische Herztätigkeit zu übernehmen. Zusätzlich wird versucht mittels
Elektroschocks (Defibrillation) die körpereigene elektrische
Erregung der Herztätigkeit zu \enquote{normalisieren}.
Beatmung, Herzdruckmassage und Defibrillation
bilden die \EE{drei Säulen der Reanimation}.}{
\begin{itemize}
  \item Kein Bewusstsein
  \item Keine normale Atmung
  \item → Kein Kreislauf
  \item → Herz-Lungen-Wiederbelebung (Reanimation)
\end{itemize}
}{}{}{}


\section{Repititorium: Anatomie und Physiologie der Atmung
und Atemwege}
\begin{itemize}
  \item Ventilation, Oxygenation, Respiration;
\end{itemize}
---
\begin{itemize}
   \item die anatomischen Strukturen der Atmung bei Kinder und
  Erwachsene benennen und beschreiben können. 
  \item die
  Atemmechanik verstehen und erklären können.
  
  \item die Physiologie der Atmung verstehen und erklären
  können.
\end{itemize}

 \index{Atemwege}
 \label{topic:atemwege-lunge}


\AuthNotes{KOCH: mag Nasenrachenraum und Mundrachenraum
streichen, evtl. wie Bild?}

\Layout{27}{}
{}
{}
{./images/anatomie/Respiratory_system_complete_de-edited.pdf}
{\AuthNotesPicLegalOk{}{PD}Die Atemwege, Übersicht.\label{fig:respiraktionstrakt-uebersicht}}
{Mariana Ruiz Villarreal, Public domain}


%\begin{multicols}{2}
Die Atemwege gliedern sich in die \EE{oberen} und die
\EE{unteren} Atemwege.

\TextSlide{Obere Atemwege}
{%
Der Weg der Einatemluft beginnt in der Nase bzw. im Mund. Die Nase ist
mit Schleimhaut ausgekleidet, knöcherne Vorsprünge sorgen für
eine Vorwärmung der Luft. In der Schleimhaut be\-finden sich die
Sinneszellen für den Geruchssinn. 

\Z{In den Schädelknochen befinden sich luft\-ge\-füllte
Hohlräume, die so\-genannten \T{Neben\-höhlen} (Stirn-,
Kiefer\-höhlen, Siebbeinzellen). Diese stehen in Verbindung mit dem Nasenrachenraum.
Entzündet sich die Schleimhaut spricht man von der
\E{Neben\-höhlen\-ent\-zündung}, welche oft langwierig
und schmerzhaft sein kann.} 

Die Luft gelangt nun in den \T{Mund-Rachen\-raum} und den
\T{Rachen}. \EE{Hier kreuzt der Weg der Nahrung mit dem
Weg der Atemluft!} Die Nahrung gelangt vom Mundraum in die
hinten gelegene Speiseröhre, die Luft möchte vom Nasen\-rachen\-raum in die vorne gelegene Luftröhre. 

Um ein Verschlucken von Nahrung in die Lunge zu verhindern
verschließt der am Kehlkopf gelegene \T{Kehldeckel}
während des Schluckaktes die Luftröhre. Dieser Reflex ist lebenswichtig! 
Im \E{Kehlkopf}  befinden sich die \T{Stimmbänder},
welche für die Lautbildung beim Sprechen wichtig sind.

\Fazit{Im Rachenraum kreuzt der Speise- und der Luftweg!}
}{%
\begin{itemize}
    \item \EE{Nasen}- und \EE{Mundhöhle} \index{Nasenhöhle} 
    \index{Mundhöhle} 
    \item \EE{Rachen} (gemeinsamer Weg von Nahrung und
    Atem!) \index{Rachen} 
    \item \EE{Kehlkopf} (\T{Larynx}) mit
    \EE{Kehldeckel} \index{Kehlkopf} \index{Larynx} 
    \index{Kehldeckel}
\end{itemize}
}{}{}{}

\TextSlide{Untere Atemwege}
{%
\label{topic:alveolen}
In der \T{Luftröhre} (\T{Trachea}) wird die vom Rachen
und Kehlkopf kommende Luft zur Lunge weitergeleitet. Die Luftröhre teilt sich in einen
\E{rechten} und einen \E{linken} \T{Hauptbronchus},
welche zu den jeweiligen Lungenflügeln ziehen. Jeder Hauptbronchus verzweigt sich immer weiter
bis schlussendlich sehr dünne \E{Bronchioli} in die
\T{Lungenbläschen} (\T{Alveolen}) münden. Hier findet
der Gasaustausch statt.

Die Abschnitte der Atemwege, in denen kein Gasaustausch
stattfindet und die nur dem Transport dienen, bezeichnet man
als {\quotedblbase}\index{Totraum}\T{Totraum}{\textquotedblleft}.
}{%
\begin{itemize}
    \item \index{Luftröhre}\EE{Luftröhre}
    (\index{Trachea}\T{Trachea})
    \item \T{Hauptbronchus} (li./re.)
    \item \IX{Bronchioli}
    \item \EE{Lungenbläschen}
    ({Alveolen})
\end{itemize}
}{}{}{}



%\end{multicols}


\subsection{Die Lunge ist für einen Teil des
Gasaustausches zuständig}
\index{Gasaustausch!Lunge}  
\index{Lunge}
\index{Lunge!-nflügel}

\TextSlide{\TxBeschreibung}
{
Die Lunge hat \EE{zwei Flügel}, welche sich jeweils in \EE{3 (rechts)} bzw.
\EE{2 (links) Lappen} gliedern: 

\begin{center}\FormatTable
\begin{tabularx}{\linewidth}[H]{XX}
\toprule\addlinespace
\multicolumn{2}{c}{\TabHeading{Die Lunge}}
\\\addlinespace\midrule\addlinespace
\multicolumn{1}{c}{\TabSub{Rechter Flügel}}
 &
\multicolumn{1}{c}{\TabSub{Linker Flügel}}
\\\addlinespace\cmidrule(lr){1-1}\cmidrule(lr){2-2}\addlinespace
\begin{itemize}
    \item \EE{3} Lappen
    \begin{itemize}
        \item Oberlappen \index{Oberlappen} 
        \item Mittellappen \index{Oberlappen} 
        \item Unterlappen \index{Unterlappen} 
    \end{itemize}
\end{itemize}
&
\begin{itemize}
    \item \EE{2} Lappen
    \begin{itemize}
        \item Oberlappen
        \item Unterlappen
    \end{itemize}
\end{itemize}
\\\addlinespace\bottomrule
\end{tabularx}
\end{center}

Die Oberfläche ist mit dem \T{Lungenfell}
\index{Lungenfell} \index{Pleura!lungenseitige}, der
\E{lungenseitigen} \T{Pleura}, überzogen. In einem Lungenflügel verzweigt
sich der Hauptbronchus in immer weitere kleinere Bronchien bis schließlich die
Bronchioli in die \E{Lungenbläschen} (\E{Alveolen}) enden.
}
{
\begin{itemize}
  \item 2 Flügel mit rechts 3 Lappen und links 2
  \item Lungenfell (Pleura)
  \item Lungenbläschen (Alveolen)
\end{itemize}
}{}{}{}

\subsubsection{Die Alveolen sind die Endstücke
des Atemwegs und sind für den Gasaustausch zuständig}

\TextSlide{\TxBeschreibung}
{
\emph{Siehe Abb.
\FullPrettyRef{fig:respiraktionstrakt-uebersicht}, rechts
oben.} In den \T{Lungenbläschen} (\T{Alveolen}) findet der
\index{Gasaustausch!Lunge}\E{Gasaustausch} zwischen der
Einatemluft und dem Blut statt. Sauerstoff gelangt dabei aus den Alveolen in das Blut und
Kohlendioxid wird im Gegenzug vom Blut an die Alveolen
abgegeben.

Dabei besteht zwischen Luft und Blut kein direkter Kontakt, die Gase
müssen Hindernisse wie die Alveolenwand und die Blutgefäßwand
überwinden. \E{Vergrößert sich dieser Weg aufgrund
einer Erkrankung (z.\,B. beim Lungenödem\VARIANT{VAR-STANDARD}{, 
\FullPrettyRef{topic:lungenoedem})}, ist die Sauerstoffaufnahme vermindert!} 

\Fazit{Lunge: Sauerstoff (\ce{O2}):   Alveolen ${\rightarrow}$ Blut}

\Fazit{Lunge: Kohlendioxid (\ce{CO2}):  Blut ${\rightarrow}$
Alveolen } }
{
\begin{itemize}
  \item Gasaustausch: Kohlendioxid gegen Sauerstoff
  \item Hindernisse: Alveolenwand, Blutgefäßwand
\end{itemize}
}{}{}{}

\clearpage
\subsubsection{Die Pleura}\label{topic:pleura}

\TODO{GABS}{2011-03-07}{Pleura: zuviel Text? Redundant? kürzen?}{}

\Layout{20}{\TxBeschreibung{} und Funktion}
{%
\index{Rippenfell}%
\index{Lungenfell}%
\index{Pleura!thoraxseitige}%
\index{Pleura!lungenseitige}%
Die Pleura hat zwei \protect\enquote{Blätter}. Eine Seite  liegt an der
Rippenseite (\T{Rippenfell}) an und ist fest mit dem Zwerchfell und der Brustwand
verwachsen. Die zweite Seite  überzieht die Lungenoberfläche (\T{Lungenfell}).
Zwischen beiden Pleurablättern befindet sich ein \E{Flüssigkeitsfilm}. Dieser
sorgt dafür, dass die beiden Blätter seitlich gleiten können. (Es gibt
keine Verwachsungen der beiden Seiten und die
Pleuraflüssigkeit fungiert sogar ein bisschen als
\protect\enquote{Schmiermittel}.)

Der Zwischenraum ist
\EE{dicht abgeschlossen}, es kommt von außen keine Luft dazu. Dadurch haften die
beiden Blätter aneinander, bei Zug (durch Unterdruck während der Einatmung) an
einem der beiden Blätter lassen sich diese nicht trennen, folglich wird das Zweite passiv
mitbewegt.\footnote{Wie bei zwei Glasplatten mit Wasser dazwischen.} 
Der Vorteil dieser Konstruktion ist, dass zwar die Lunge jeder Bewegung des
Brustkorbs folgt, jedoch nicht fixiert ist, sondern verschieblich bleibt. 

\Z{Beim Gesunden gibt es keine Verwachsungen zwischen diesen beiden
Pleurablättern. Bei Entzündungen kann es jedoch dazu kommen. Die Pleura ist
außerdem die einzige Struktur der Lunge die schmerzempfindlich ist.} }
{
\begin{itemize}
    \item \E{Thoraxseitige} Pleura (Rippenfell)
    \item \E{Lungenseitige} Pleura (Lungenfell)
    \item Zwischenraum luftdicht, mit Schmierflüssigkeit
    \item Brustkorb mit Rippenfell \protect\enquote{zieht an} → Lungenfell mit Lunge folgt
    Bewegung     
\end{itemize}
}
% {./images/noauth/AASdev20090228-img41.png}
{./images/hirtler-aass/Pleura.pdf} {Pleura, Schema.} {Hirtler.}


\EE{Frage}: \Speech{Was würde passieren wenn der
Zwischenraum zwischen den Pleurablättern aufgrund einer Verletzung (z.\,B. ein
Loch aufgrund einer Stichwunde im Brustbereich) nicht mehr luftdicht
ist?} Antwort: s. Fußzeile (\footnote{Antwort: Die Lunge
würde in sich zusammenfallen, da sie ja nicht mehr
aufgespannt wird: \I[Pneumothorax!anatomische Erklärung]{Pneumothorax}\VARIANT{VAR-STANDARD}{
(\FullPrettyRef{topic:pneumothorax})}.})


% \TextSlide{Welche Funktion hat aber nun die Pleura?}
% {%
% Da die eine Seite mit dem Brustkorb, die andere Seite mit
% der Lunge verwachsen ist und der Raum zwischen den beiden
% Seiten luftdicht abgeschlossen ist, haften die beiden Seiten
% aneinander, aber sind gegeneinander verschieblich. (Es gibt
% keine Verwachsungen der beiden Seiten und die
% Pleuraflüssigkeit fungiert sogar ein bisschen als
% \protect\enquote{Schmiermittel}.)
% 
% Wenn sich nun der Brustkorb oder das Zwerchfell bei der
% Atmung bewegt, bewegt sich natürlich das Rippenfell mit (es
% ist ja mit dem Brustkorb und dem Zwerchfell verwachsen). 
% }
% {
% \begin{itemize}
%   \item Leine Luft zwischen beiden Pleuren
%   \item Pleuraflüssigkeit = Schmiermittel
%   \item Unterdruck bei Atmung
%   \item bei Verletzung bzw. Lufteindringen in den Pleuraspalt Pneumothorax
% \end{itemize}
% }{}{}{}

\Fazit{Wenn
der Brustkorb mit dem daran verwachsenen Rippenfell \protect\enquote{anzieht}, folgt das
Lungenfell mit der Lunge der Bewegung \Sign{→} die Lunge wird wie eine Ziehharmonika
aufgezogen und Luft angesaugt.}

\subsection{Atemmechanik}
\label{topic:atemmechanik} 
\index{Atemmechanik} 
\index{Zwischenrippenmuskulatur}

Wenn man sich den Thorax annähernd als Zylinder vorstellt,
so kann eine Volumenvergrößerung (Einatmung) durch eine
Vergrößerung des Querschnitts (Heben der
Zwischenrippenmuskulatur)
oder Vergrößerung der Höhe (Anspannung des Zwerchfells) erfolgen.

\begin{equation}
V = A \times h
\end{equation}
\begin{equation}
\mbox{Volumen} = \mbox{Fläche} \times \mbox{Höhe}
\end{equation}

\AuthNotes{KOCH: entweder wird das genauer erklärt wie die
Funktion der Pleura oder es wird ganz einfach mündlich
erklärt.

GABS: Sollte schon schriftlich erklärt werden \ldots wir
wollen ja auch immer dass sie die \protect\enquote{Atemhilfsmuskulatur}
etc. unterstützen, dann sollten sie wenigstens auch
nachlesen können \protect\enquote{wie man richtig atmet}

KOCH: soll erklärt werden.}

\begin{figure}[H]
\begin{center}
% \includegraphics[width=\linewidth]{./images/noauth/AASdev20090228-img42.png}
\includegraphics[width=\linewidth]{./images/hirtler-aass/Atmung.pdf}

\Captionof{figure}{Atemmechanik: Bei der Einatmung senkt
sich das Zwerchfell und der Brustkorb hebt sich. Das
mittlere Bild zeigt die Zwischenrippenmuskulatur.
\SourcePicture{Hirtler.}}
\TODO{GABS}{0.14.0}{Atemmechanik: Bilder trennen und
getrennt beschreiben.}{}
\end{center}
\end{figure}


\TextSlide{\TxBeschreibung}
{
Für die Atemmechanik ist ein gutes Zusammenspiel von Skelett, Muskeln, Gewebe
und Pleura(spalt) wichtig. 

\Fazit{Ein bis zum Hals Verschütteter erstickt.}

\subparagraph{Atemmuskulatur} Der wichtigste Atemmuskel ist
das \EE{Zwerchfell}\index{Zwerchfell}. In Ruhe erledigt es
fast die gesamte Atemarbeit. Bei verstärkter Atmung kommen die
Zwischenrippenmuskeln zum Einsatz \cite{LippertAnatomie7}.

\Fazit{Der wichtigste Atemmuskel ist das \EE{Zwerchfell}!}

\subparagraph{Atemhilfsmuskulatur}% 
\label{topic:atemhilfsmuskulatur}%
\index{Atemhilfsmuskulatur}%
Neben den Hauptatemmuskeln (Zwerchfell, Zwischenrippenmuskulatur) gibt es noch
eine Reihe weiterer Atemhilfsmuskeln die die Atmung unterstützen. Diese
verbinden den Thorax und die Oberarme. Fixiert man die Arme, z.\,B. durch
Aufstützen, können die Muskeln die Hebung des Brustkorbes unterstützen. 
\VARIANT{VAR-STANDARD,VAR-ERNAEHRUNGSPAED}{Vgl.
dazu }\VARIANT{VAR-STANDARD,VAR-ERNAEHUNGSPAED}{\emph{Asthma}, \FullPrettyRef{topic:asthma-bronchiale}; }\VARIANT{VAR-STANDARD}{\emph{COPD},
\FullPrettyRef{topic:copd}}\VARIANT{VAR-STANDARD,VAR-ERNAEHRUNGSPAED}{.}

\subparagraph{Pleura} Die Funktion der Pleura ist unter \ref{topic:pleura}
erklärt. \Z{Der Unterdruck in der Pleura entsteht dadurch,
dass sich einerseits die (elastische) Lunge zusammenziehen
möchte, während andererseits der Brustkorb einen Gegenzug
ausübt. Diese beiden Kräfte stehen bei entspannter
Atemmuskulatur im Gleichgewicht.}


% \Z{Thoraxwand und Lunge besitzen elastische Eigenschaften: Sie werden gedehnt
% und entwickeln elastische Rückstellkräfte. An der Lunge können zwei Mechanismen ein
% Zusammenfallen des Lungengewebes bewirken. Erstens sind im Lungenzwischengewebe
% reichlich elastische Fasern eingelagert, die die Eigenelastizität der Lunge
% bewirken. Durch die Koppelung von Lunge und
% Thoraxwand am Pleuraspalt entsteht ein elastisches Gesamtsystem. 
% 
% Bei entspannter
% Atemmuskulatur und offenem Kontakt mit der Umgebungsluft wird sich ein
% Gleichgewicht der elastischen Kräfte einstellen; der Brustkorb befindet sich in
% der Atemruhelage. Durch den Zug der Lunge, die sich verkleinern möchte, und den
% Gegenzug der Thoraxwand entsteht ein ständiger Unterdruck im Pleuraspalt.}


\subparagraph{Atemzyklus}
Ein Atemzyklus besteht aus der Einatemphase und der
Ausatemphase. Dabei kommt es zu einem
Ungleichgewicht zwischen der nach außen gerichteten Zugkraft des Thorax und
der nach innen gerichteten Zugkraft der Lunge. Dieses
Ungleichgewicht wird beim Einatmen zunächst von einer
zusätzlichen Zugkraft des Zwerchfells und der äußeren Zwischenrippenmuskulatur hervor
gerufen und anschließend beim Ausatmen durch Entspannung der
Muskeln wieder abgebaut.

Während des
Einatmens wird durch aktive Bewegung des Zwerchfells und der
äußeren Zwischenrippenmuskulatur das Volumen des Brustkorbs
und somit auch das der Lunge vergrößert. Die in der Lunge
befindliche Luft hat nun deutlich mehr Platz, weswegen ihr Druck jetzt
kleiner ist, als der Luftdruck in der Umgebung
(\EE{Unterdruck}). Dieser Druckunterschied wird durch
einströmende Luft ausgeglichen -- man hat eingeatmet. 
Da für das Einatmen Muskelarbeit
verrichtet werden muss, nennt man das Einatmen auch
\E{aktive Phase} der Atmung.

Das Ausatmen erfolgt passiv (\E{passive Phase} der Atmung).
Das Zwerchfell und die äußere Zwischenrippenmuskulatur
entspannen wodurch sich das Thorax- bzw.
Lungenvolumen auf das ursprüngliche Volumen
verkleinert.\footnote{Man denke an eine Feder, welche
nachdem sie durch eine Zugkraft gedehnt wurde, nach dem
Loslassen durch ihre eigene (elastische) Rückstellkraft
wieder von selbst in die Ausgangslage zurück kehrt.} Durch die Volumenverminderung steigt der Druck in
der Lunge und die Luft strömt nach außen bis die
Druckunterschiede zwischen Lunge und Umgebung wieder
ausgeglichen sind. Die Luft wird beim Ausatmen sozusagen
wieder hinaus gedrückt.

Die Anzahl der Atemzyklen pro Minute nennt man Atemfrequenz
(AF). Sie liegt beim Erwachsenen, der sich in Ruhe befindet,
bei etwa 12 Atemzyklen pro Minute.

% Volumenschwankungen des Brustraums bewirken
% gleichzeitig auch Volumenschwankungen im Inneren der Lunge. Eine Vergrößerung
% des Thoraxinnenraums wird durch Kontraktion der Atemmuskulatur bewirkt. Dieser
% Vorgang wird als aktive Phase der Atmung bezeichnet. Dabei kommt es vor allem
% zu einer Abflachung und einem Tiefertreten des Zwerchfells, aber auch zu einer
% Anhebung der Rippen durch die äußeren Zwischenrippenmuskeln. Als Folge entsteht
% in der Lunge im Vergleich zur Umgebungsluft ein Unterdruck, der durch
% Einströmen von Luft ausgeglichen wird. Thoraxwand und Lunge werden dabei
% gedehnt.
% 
% Am Ende einer ruhigen Einatmung
% erschlafft die Einatmungsmuskulatur. Die elastischen Rückstellkräfte führen das
% Thoraxvolumen wieder in die Ausgangslage zurück. Dabei kommt es zu einer
% Druckerhöhung in der Lunge, die eine Luftbewegung hin zur Umgebungsluft bewirkt.
% Diese erfolgt in aller Regel passiv, d. h. ohne aktive
% muskuläre Anstrengung. Dieser Vorgang erfolgt in Ruhe bei einem Erwachsenen
% etwa zwölfmal in der Minute und wird als Atemfrequenz bezeichnet.
}
{
\begin{itemize}
  \item Atemmuskulatur: Zwerchfell, Zwischenrippenmuskeln
  \item Atemhilfsmuskulatur
  \item Elastische Eigenschaften von Thorax und Lunge
  \item Unterdruck für Atmung wichtig
\end{itemize}
}{}{}{}

\VARIANT{VAR-STANDARD}{\subsection{Atemfrequenz, -tiefe und -volumen}

Die Atemfrequenz, die Atemtiefe und das Atemvolumen
werden  unter
\FullPrettyRef{topic:atemvolumina} besprochen.
}


% \clearpage
\subsection{Atemgase}

\Definition{Als Atemgas bezeichnet man jenes Gasgemisch,
welches eingeatmet wird.}

\TextSlide{\TxBeschreibung}
{%
\index{Stickstoff}%
\index{Atemgas}%
\index{Sauerstoff!Raumluftgemisch}%
\index{Kohlendioxid!Raumluftgemisch}%
Normalerweise ist das ein Gasgemisch aus Stickstoff,
Sauerstoff, Kohlendioxid und Edelgasen (\protect\enquote{Raumluft}).
Zwischen der normalen Einatem- und Ausatemluft besteht besonders hinsichtlich
des Sauerstoff- und Kohlendioxid-Anteils ein Unterschied, siehe 
\prettyref{tab:luft-gas}.

\begin{table}[H]
\FormatTable
\Captionof{table}{Gasgehalt in der Raum- und
Ausatemluft.\label{tab:luft-gas}}

\begin{tabular}[H]{ccc}
\addlinespace\toprule\addlinespace
\centering \TabHeading{Raumluft} & & \TabHeading{Ausatemluft}
\\\addlinespace\midrule\addlinespace
\EE{21\,\PC*}  & Sauerstoff (\ce{O2}) & \EE{16\,\PC*}
\\\addlinespace
78\,\PC* & Stickstoff (\ce{N}) & 78\,\PC*
\\\addlinespace
\EE{0,03\,\PC*} & Kohlendioxid (\ce{CO2}) & \EE{4\,\PC*} 
\\\addlinespace
1\,\PC* & Edelgase  & 1\,\PC*
\\\addlinespace\bottomrule
\end{tabular}
\end{table}
}
{
\begin{itemize}
  \item Raumluft
  \item Atemgase: Sauerstoff und Kohlendioxid
\end{itemize}
}{}{}{}







\section{Supraglottische Atemwegssicherungen}

Supraglottische Methoden zur Atemwegssicherung werden oft
als \enquote*{alternativer Atemweg} bezeichnet, da sie oft
eine Alternative zur endotrachealen Intubation (\I{ETI})
angesehen werden. Die ETI gilt weithin als Goldstandard der
Atemwegssicherung, dennoch hat sie einige gewichtige
Nachteile und ist in vielen Situationen \E{nicht die
Methode der Wahl}. 





\subsection{Larynxtubus{\texttrademark}}



\subsection{Larynxmaske}
\subsection{Combitubus}

\section{Algorithmen}
\subsection{ASBÖ-Algorithmus \enquote{Intubation}}
---
\begin{itemize}
  \item die Technik der endotrachealen Intubation mit ihren Vor-
  und Nachteilen kennen und beschreiben können.
  
  \item bei der Vorbereitung und Durchführung der Intubation
  einem Arzt oder NFS-NKI assistieren können. \item die
  alternativen Atemwege kennen und bei der Anwendung
  assistieren können.
  
  \item die Indikation, Kontraindikation sowie die Komplikationen
  alternativer Atemwege kennen und beschreiben können.
\end{itemize}


\ChapterEnd
